\chapter{Random}
\label{ch:random}

\newthought{Uniform draws, exactly the budget you asked for.} \sampler{Random} fills
u-space with independent, identically-distributed draws, mapped through each variable's
declared distribution. It is the simplest method in the book and the recommended first
step for any new card: it exercises the full execution path (workflow, likelihood,
archiving) without any sampler-specific behavior to reason about.

\section{Configuration}
\label{sec:random-config}

\begin{yamlwin}
Sampling:
  Method: Random
  "Point number": 500          # REQUIRED -- note the space in the key
  Seed: 11
  Variables:
    - name: x
      distribution: {type: Flat, parameters: {min: 0.0, max: 1.0}}
    - name: y
      distribution: {type: Flat, parameters: {min: 0.0, max: 1.0}}
\end{yamlwin}

\begin{keylist}
  \item["Point number"] \req\ (\code{ValueError} if absent); alias
        \yamlkey{point\_number}. The number of draw \emph{attempts} --- the
        \emph{nominal} budget. A \yamlkey{selection} cut only reduces how many of those
        attempts are accepted; it does not top the budget back up
        (Section~\ref{sec:selection}).
  \item[Seed] \code{0} (the default) means unseeded --- a fresh \textsc{numpy} stream
        every run, so identical cards produce different points on different launches.
        Set it to any nonzero integer for anything you may want to resume or reproduce.
\end{keylist}

When to use it: first contact with a new space, workflow debugging, an unbiased
baseline map to compare other methods against, and any time you want the attempted
sample count to be exactly what you asked for rather than emergent.

\section{How it behaves in the runtime}
\label{sec:random-runtime}

Draws happen lazily, one at a time, in the control process: each call to the sampler's
internal \code{propose\_next} draws a fresh u-space vector from
\code{numpy.random}, seeded once at startup if \yamlkey{Seed} is nonzero. A drawn point
that fails \yamlkey{selection} is discarded and the loop draws again --- but the attempt
still counts against \yamlkey{"Point number"}, which is why a tight cut yields fewer
\emph{accepted} samples than the nominal budget (the same accounting rule as every
other fixed-set method, Section~\ref{sec:selection}).

Each accepted point's \textsc{uuid} is deterministic in \code{(seed, accepted\_index)},
and accepted points are submitted to Redis in the usual
\yamlkey{EnvReqs.V2.batch\_size} groups. The sampler's checkpoint carries both the
attempt/accept counters \emph{and} \textsc{numpy}'s own random-generator state, so a
resumed run continues drawing from the exact same stream rather than restarting it ---
the sequence of points after resume is identical to an uninterrupted run with the same
\yamlkey{Seed} (Chapter~\ref{ch:checkpoint}).

\begin{jtip}
Because the draw stream is fully determined by \yamlkey{Seed}, two cards that differ
only in \yamlkey{selection} or in the workflow draw the \emph{same} underlying u-space
points --- useful when comparing how a likelihood or a cut reshapes an otherwise
identical sample.
\end{jtip}

\section{Choosing \sampler{Random} over the alternatives}

\begin{itemize}
  \item \textbf{vs \sampler{Bridson}:} at equal budget, \sampler{Random} clumps and
        leaves voids that \sampler{Bridson}'s minimum-distance packing avoids ---
        prefer \sampler{Bridson} once you are past the debugging stage and about to
        interpolate or plot the result (Chapter~\ref{ch:bridson}).
  \item \textbf{vs \sampler{Grid}:} \sampler{Random} scales to any dimension without
        the combinatorial blow-up of a Cartesian product, at the cost of even coverage
        (Chapter~\ref{ch:grid}).
  \item \textbf{vs \sampler{AdaptiveBridson}:} \sampler{Random} spends its budget
        uniformly; if what you actually want is dense coverage \emph{near a contour},
        \sampler{AdaptiveBridson} spends the same budget far more efficiently
        (Chapter~\ref{ch:adaptive-levelset}).
\end{itemize}

\section{Summary}
\label{sec:random-summary}

\sampler{Random} is the fixed-set baseline: a nominal draw budget, no geometry to
configure, resume that replays its exact random stream. Reach for it first on any new
card, and keep it as the reference map to compare every other method against.

\begin{table}[ht]
  \footnotesize
  \begin{tabular}{@{}p{0.32\linewidth}p{0.62\linewidth}@{}}
    \toprule
    \textbf{Field} & \textbf{\sampler{Random}} \\
    \midrule
    Kind & fixed set \\
    Point count & \yamlkey{"Point number"} (nominal draw attempts) \\
    Dimensions & any \\
    Required keys & \yamlkey{"Point number"} \\
    u$\to$x mapping & yes \\
    \yamlkey{Seed} role & draw stream (\code{0} = unseeded) \\
    \yamlkey{selection} & consumes attempts, no top-up \\
    Extra output & --- \\
    \bottomrule
  \end{tabular}
  \caption{\sampler{Random} at a glance.}
  \label{tab:random-summary}
\end{table}
